\section{Related work}
\label{sec:related}

Prior work relevant to the problem of Section~\ref{sec:intro} can be
grouped into five categories: network-level measurement tools, greedy
instance-reconstruction methods, joint neighbourhood-optimisation methods,
complete micrograph-to-measurement workflows, and monocular depth-ordering
methods.

\subsection{Network-level methods without instance reconstruction}
\label{sec:rel-tools}

Existing tools answer part of the problem. Semantic segmentation is mature,
since U-Net and
its self-configuring descendants reliably separate strand foreground from
background \citep{ronneberger2015unet,isensee2021nnunet}, in materials
imaging as much as anywhere
\citep{stan2020optimizing,davydzenka2022deep}, and nanotube micrographs
have their own trained segmenters
\citep{trujillo2017segmentation,nguyen2023cntsegnet,dedonato2025deep}. But a
foreground mask encodes no identity, so it inherits both failures rather than
resolving them. Instance methods assume the objects are compact, in materials microscopy
\citep{he2026instance} as in the general-purpose families: bounded
blobs, centre-directed flows, star-convex outlines, none of which describes a
long open curve that may overlap another \citep{stringer2021cellpose,%
schmidt2018stardist,kirillov2023sam}. Ridge and vesselness filters sharpen
curvilinear signal but assign nothing across a junction
\citep{steger1998unbiased,frangi1998vesselness,sato1998linefilter}, and the
established fibre-analysis packages return network statistics rather than
resolved objects \citep{stein2008fire,bredfeldt2014ctfire,xu2015soax,%
flormann2021finta,hotaling2015diameterj}. Where the imaging is volumetric the
problem is different in kind, because strands that appear to meet occupy
disjoint volumes: in X-ray tomography of fibrous solids and of twisted ropes,
individual fibres are separated by seeded watershed and morphological
reasoning on the reconstructed volume
\citep{depriester2022fibre,allan2025rope}. A projection has no such recourse,
since two strands that cross genuinely share the same pixels.

\subsection{Greedy continuation from a mask}
\label{sec:rel-greedy}

Recovering strands through crossings has been addressed repeatedly. Existing
methods differ primarily in the evidence used for continuation and in whether
earlier pairings can subsequently be revised.

A large class of methods follows a common procedure: skeletonise the mask, cut
the crossing out,
and pair the arms left behind by how nearly each continues the direction of
another. Each pairing is accepted as soon as it is scored. SIFNE reads its tip
directions from the excised stubs and recovers dense microtubule networks from
the mask alone \citep{zhang2017sifne}. Wegmayr et al.\ pair skeleton branches
by lowest mutual angle \citep{wegmayr2020microplastic}. GraFT covers a
Frangi-filtered skeleton greedily, longest path first
\citep{osterlund2026graft}. This is inexpensive, and where crossings are rare
it suffices. Its weakness is structural: once an arm is paired, a later arm
that contradicts the choice cannot reopen it.

Two remedies have been tried. The first is to bring in more evidence
than a mask contains. The stronger of Wegmayr et al.'s two variants drops the
angle rule and
learns pixel embeddings on patches around each junction; it reads the greyscale
image, and it must be trained on crops in which every fibre is already labelled
individually. Even then it recovers tangled fibres at about half the precision
of isolated ones, in data where two thirds of fibres cross nothing at all.
Methods that pair ends with an orientation network, detect fibres with a
region-based detector, or trace them recurrently
\citep{liu2019intersection,frei2020fiber,liu2023picktrace} need the same
labelled training data. Reading a mask and nothing else avoids that cost, a
trade-off adopted in this paper. DNAi keeps to the mask and searches more
widely,
enumerating every pairing a junction allows from its endpoint angles; it does
so exhaustively at degree three and four, and not at all above
\citep{playout2026dnai}. The closest single prior method to PLECTA is that of
Liu et al., which likewise reads a mask, reconnects across gaps and permits
overlapping output, but pairs termini greedily and depends on a trained
orientation network \citep{liu2019intersection}.

\subsection{Joint neighbourhood optimisation}
\label{sec:rel-joint}

A second remedy is to stop deciding one pairing at a time. Basu et al.\ and
DeFiNe both weigh a whole neighbourhood at once, which is what a dense node
requires. Basu et al.\ delete every branch point from a spanning tree over
centreline points, then reassemble the pieces with an integer program that
prices each join by how sharply it bends, repeating until the set of pieces
settles; a join that a later fit contradicts can therefore still be undone
\citep{basu2015extraction}. DeFiNe covers an already-built weighted graph with
the smoothest paths it can find, and lets two paths share an edge, so a
crossing pixel can belong to both strands instead of being awarded to one
\citep{breuer2015define}. Neither, however, starts from an axis mask: Basu et al.\ require image
intensities, whereas DeFiNe requires a graph produced by an upstream stage. Both
also reassemble at most three pieces at a time, which limits how much of a
shattered neighbourhood a single round can repair. Deciding broadly and
starting from a mask have so far been alternatives.

One method combines both properties, and comes closest to the present
approach. De
et al.\ reason globally and do start from a mask. Branches and junction
contacts become a graph, and
labels spread across it so that a decision at one node constrains its
neighbours rather than resting on that node's own angles \citep{de2016graph}.
The labels have to start somewhere, however. In their images they start at
anatomical roots, the stained cell bodies of neurons or the optic disc of a
retina, and the method returns as many strands as there are roots. A tangle of
open-ended strands offers no such starting point, and how many strands it
holds is precisely the unknown. Reasoning globally from a mask alone, with no
seed and no strand count supplied in advance, is the problem addressed in this
paper. Four of these methods are measured against PLECTA directly in
Section~\ref{sec:res-comparators}, on identical masks and through one scorer.
Skeleton graphs, junction excision, continuation angles, distance gates,
combinatorial optimisation, iterative refitting and overlapping output are
therefore all established; none is claimed as new in this paper.

\subsection{Complete workflows, and depth from a single view}
\label{sec:rel-workflows}

The closest complete workflow is that of Peters et al., who join a U-Net-like
segmentation trained on over a thousand annotated SEM micrographs to
algorithms that trace fibres through intersections, and validate against human
evaluators on fibre number, length and width \citep{peters2026fibre}. Both
workflows separate semantic segmentation from instance reconstruction but
evaluate the reconstruction stage differently: Peters et al.\ assess the
replacement of human counting and measurement, whereas PLECTA is isolated
behind a binary mask and evaluated on the preservation of strand identity
across crossings, against overlap-aware references in which every pairing is
known.

Which strand lies on top at a crossing is a further question, and for
nanostructured networks it is normally answered destructively, by
focused-ion-beam sectioning and tomographic reconstruction
\citep{gabbett2024fibsem}, because a micrograph's intensities record surface
topography rather than height. Monocular depth ordering has long derived
pairwise depth from local occlusion evidence, assembled those relations into a
depth-order graph and corrected the contradictions among them
\citep{palou2013monocular,lee2022instancewise}, and single-view hair modelling
resolves the layering of intersecting strands by analogous reasoning
\citep{chai2012singleview}. Those methods read the shape of the outlines where
two objects overlap, which a bundle a few pixels across in a noisy micrograph
does not reliably provide; the evidence used in this paper is instead the
continuity of each strand's own brightness across the crossing.
